% Abstract. Every number here comes from verify/registry.json.
\chapter*{Abstract}
\phantomsection
\addcontentsline{toc}{chapter}{Abstract}
\begingroup\setstretch{1.15}

Something in the clouds of Venus absorbs near-ultraviolet sunlight and is responsible for
about half of the solar heating at the cloud top, yet nearly a century after its markings
were first photographed, the substance is still unknown. This project asks not which
molecule has the right spectrum but whether any proposed carrier can be present in
the right amount, at the right altitude and with the right spectral shape, all at once,
inside a model that is free to move it around.

The work uses VULCAN-Venus, a one-dimensional model of the photochemistry of the Venus
atmosphere. Phase~1 found that the model's absorber term was present but switched off,
switched it on under a single amplitude, and ran fourteen experiments on the model and
on published laboratory data. Every number is
recomputed from the saved model output by one script into a registry of 10830 values,
and a second script checks every printed number precise enough to test against it.

The absorber's effect on the radiation is proportional to its amount and is concentrated
in one band of altitude, 69 to 93~km, which holds 91.370\% of the response. Where an
absorber sits matters more than its spectrum: moved up, the same absorber changes the
radiation 7.34 times more than moved down, against a factor of 4.95 between the
spectra tested at a fixed height. When the model's own chemistry is allowed to place
ClS\sub{2}, the gas closest to the needed amount, its centre falls 51.521~km lower and
it becomes radiatively the same as no absorber, to 0.001487\%. A screen of all 23
species the model can represent finds none that passes amount, placement and spectrum
together. From published data, the proposed iron and sulfur minerals cannot absorb enough at 375~nm even as pure solids, and a
published rejection of ferric chloride, based on a measurement in water, does not hold
if it absorbs in the droplets as it does in hydrochloric acid, where the amount needed
falls from 71.0\% of the droplet mass to between 2.82\% and 17.05\%. Changes in
abundance smaller than 2.5278\% are shown to be numerical, not physical, and are
reported only as upper bounds.

Phase~2 proposes to test these verdicts under the much stronger vertical mixing now
estimated for the Venus clouds, to ask whether upwelling can carry a candidate into the
response band, and to place ferric chloride by the same altitude test.

\textbf{Keywords:} Venus, unknown ultraviolet absorber, photochemistry, VULCAN,
radiative transfer, ferric chloride, reproducibility
\par\endgroup
